\section{Analysis}

\subsection{The headline: nobody splits a pair below \texorpdfstring{$d/\lam \approx 10^{-4}$}{d/lam ~ 1e-4}}
\label{sec:headline}

Family B is the pair ladder. The table gives the \emph{oracle} cost --- the cheapest tested
setting that succeeded --- at each separation, at the tightest $\pr$ for which the pair must
still be \emph{split} rather than merged. A dash is failure at every setting within the
1500-evaluation budget.

\begin{center}
\begin{tabular}{rrrrrrrr}
\toprule
$d$ & $d/\lam$ & \texttt{uniform} & \texttt{lappy} & \texttt{lappy\_rep} & \texttt{lipschitz}
& \texttt{swing\_cells} & \texttt{swing\_direct} \\
\midrule
5.8e-01 & 1.2e-02 & 446 & 302 & 1316 & 293 & \textbf{131} & 152 \\
1.8e-01 & 3.6e-03 & --  & 1016 & 1103 & 1298 & \textbf{528} & 688 \\
5.9e-02 & 1.2e-03 & --  & -- & 922 & 1308 & \textbf{184} & 225 \\
1.8e-02 & 3.6e-04 & --  & -- & \textbf{830} & -- & -- & -- \\
5.9e-03 & 1.2e-04 & --  & -- & 1426 & -- & \textbf{658} & 844 \\
1.8e-03 & 3.6e-05 & --  & -- & -- & -- & -- & -- \\
\multicolumn{2}{r}{$\dots$ down to} & -- & -- & -- & -- & -- & -- \\
5.9e-12 & 1.2e-13 & --  & -- & -- & -- & -- & -- \\
1.8e-12 & 3.6e-14 & --  & -- & -- & -- & -- & -- \\
\bottomrule
\end{tabular}
\end{center}

\noindent
\textbf{Every method's splitting limit lies between $d/\lam = 1.2\mathrm{e}{-4}$ and
$3.6\mathrm{e}{-5}$} at this budget --- four to five orders coarser than the pencil's own
resolution limit of $\varepsilon/c \approx 4\mathrm{e}{-14}$ absolute, about $10^{-15}$
relative. \textbf{The gap between what the objective supports and what the searches extract is
roughly ten orders of magnitude}, and it is the most important number here.

Below the limit the methods do not fail loudly: they merge the pair and report success at any
$\pr$ loose enough to permit it. At $d = 1.8\mathrm{e}{-3}$ and $\pr = 10^{-4}$, where merging
is legal, \texttt{uniform\_grid} passes at 86 evaluations; at $\pr = 10^{-8}$ on the same case,
where it is not, every method fails.

\subsection{Family B is not the whole story, and reading it alone gives the wrong verdict}

A first draft of this report concluded from the ladder above that \texttt{lipschitz} should be
cut. The per-family breakdown contradicts that, and the contradiction is recorded here rather
than quietly fixed. Oracle\,/\,autonomous passes at $\pr = 10^{-8}$:

\begin{center}
\begin{tabular}{lrrrrr}
\toprule
method & A (1) & B (12) & C (8) & D (3) & E (1) \\
\midrule
\texttt{uniform\_grid}   & 1/0 & 3/2 & 3/2 & 3/3 & 0/0 \\
\texttt{lappy}           & 0/0 & 4/3 & 4/2 & 3/3 & 0/0 \\
\texttt{lappy\_repaired} & 1/0 & \textbf{6}/0 & \textbf{6}/2 & 3/3 & 0/0 \\
\texttt{lipschitz}       & 1/0 & 5/\textbf{5} & 5/\textbf{4} & 3/3 & 0/0 \\
\texttt{swing\_cells}    & 0/0 & \textbf{6}/\textbf{5} & 3/3 & 3/3 & 0/0 \\
\texttt{swing\_direct}   & 1/\textbf{1} & \textbf{6}/\textbf{5} & 3/3 & 3/3 & 0/0 \\
\bottomrule
\end{tabular}
\end{center}

Three things the pair ladder alone concealed:

\textbf{The swing methods are notably worse on clusters than on pairs} --- 3 of 8 on family C
against 6 of 12 on family B, where \texttt{lappy\_repaired} reaches 6 of 8 and
\texttt{lipschitz} 5. A cell decomposition cut at eigenvector crossings is well suited to
splitting two eigenvalues and less so to a tight group of three or four with unequal internal
gaps.

\textbf{\texttt{lappy\_repaired} has the best oracle record and the worst autonomous one} ---
6 of 12 on family B at oracle, and \textbf{0 of 12} autonomously. Essentially all of its
apparent capability is per-case tuning. That is the single largest oracle--autonomous gap in
the study, and it is exactly the failure mode the two-column protocol exists to expose: a
method that looks strongest in a tuned comparison and would ship as the weakest.

\textbf{\texttt{lipschitz} is the most consistent autonomous performer}, 5/5 on B and 4 of 5
on C --- its oracle and autonomous columns nearly coincide, meaning one fixed setting works
across cases. Its known theoretical weakness, that a pilot-estimated Lipschitz constant can
underestimate the true slope and render the completeness certificate false, did not bite on
rectangles. That is a real caveat for corner domains, not a measured failure here.

\textbf{Nobody passes family E} --- the 29-eigenvalue mixed window --- at any $\pr$ or any
setting. A window holding isolated levels, six pairs and a triple defeats every method in the
study within budget.

\subsection{Cost scales with \texorpdfstring{$\pr$}{prec} even when the requirement does not}

On most cases the purity requirement is \emph{identical} at $\pr = 10^{-8}$, $10^{-11}$ and
$10^{-14}$: the smallest true gap is far above $\pr\!\cdot\!\lam$ at all three, so nothing about
what must be resolved changes. Yet cost rises monotonically. On $d = 5.8\mathrm{e}{-1}$,
\texttt{swing\_cells} needs 131, 170, 200 and 230 evaluations at the four precisions --- a
$1.8\times$ increase for ten orders of $\pr$, buying nothing.

The cause is that every method bisects to $\pr\!\cdot\!\lam$ regardless of the local gap.
Refining a cell boundary to $10^{-13}$ when the nearest other eigenvalue is $0.58$ away is pure
waste. This violates the ``not more work than asked'' clause of the operational definition, and
it is a defect of the whole field rather than of any one method.

That this is cost and not capability was confirmed directly rather than inferred. On
$d = 5.9\mathrm{e}{-3}$ at $\pr = 10^{-14}$, where the requirement is identical to
$\pr = 10^{-11}$, \texttt{swing\_cells} fails at the 1500-evaluation cap and \textbf{succeeds at
1549 evaluations} once the cap is raised --- against the 1279 it needs at $10^{-11}$. A
21\,\% cost increase, and the only thing separating ``fails'' from ``succeeds'' is the budget
line.

\subsection{The controls validate the metrics}

\texttt{single\_bracket} returns $[(a,b)]$ and scores coverage $1.0$ with zero spurious on every
case. It is rejected by \textbf{purity alone}, on 18 of 21 cases. The three it passes are the
family-D exact-degeneracy cases, which contain a single level --- there one bracket \emph{is}
the correct answer, so the metric accepting it is right rather than broken.

\texttt{fine\_grid} exceeds the evaluation cap on every case and is rejected on \textbf{cost}.
Neither degenerate strategy survives, and each is rejected by the criterion intended for it.

\section{Filter verdict}

Capability claims within the measured budget of 1500 evaluations, on rectangles, for the
bracketing phase only.

\begin{description}
\item[\texttt{uniform\_grid} --- CUT as a candidate.] Weakest or joint-weakest in every family,
splits no pair tighter than $d/\lam = 1.2\mathrm{e}{-2}$. It should be retained as the control
that every proposal must beat, which is the role it already has.

\item[\texttt{lappy} (original) --- CUT.] Dominated by \texttt{lappy\_repaired} in every family
at both oracle and autonomous, and splits nothing below $d/\lam = 3.6\mathrm{e}{-3}$. Its
$\sig_2$-value recursion gate was already known to test the grid rather than the pair; this
measures the consequence.

\item[\texttt{lappy\_repaired} --- KEEP, with the gap as the priority.] Best oracle record of
any method (6/12 on B, 6/8 on C) and the worst autonomous (0/12 on B). Nothing about its
subdivision is wrong; what fails is its self-configuration. Closing that gap is worth more than
any change to its search strategy.

\item[\texttt{lipschitz} --- KEEP.] The most consistent autonomous performer, with oracle and
autonomous nearly coinciding. This reverses the draft verdict, which cut it on the pair ladder
alone. Its unsound-certificate risk is real but was not exercised here, and should be tested on
a corner domain before it is either trusted or discarded.

\item[\texttt{swing\_cells} --- KEEP.] Cheapest method wherever it succeeds, often by
$3$--$7\times$, and reaches the ladder's limit. Its weaknesses are clusters of three and four
(3 of 8) and the isolated-eigenvalue case, which it fails at every $\pr$.

\item[\texttt{swing\_direct} --- KEEP.] Costs 15--25\,\% more than \texttt{swing\_cells} and
buys no pair-splitting capability, but it is what rescues the isolated case (1/1 against 0/0).
The two mechanisms are orthogonal rather than ranked, and the ablation says so: cell-cutting
handles pairs, \textsc{direct} handles isolated wells.
\end{description}

\noindent
No method is cut on the strength of family B alone, and no method survives family E.

\section{What this cannot conclude}

\begin{itemize}
\item \textbf{A correction.} An earlier draft of this report claimed that
$\pr = 10^{-14}$ could never be exercised as a requirement on this instrument, on the grounds
that the only separation distinguishing it from $10^{-11}$ --- a pair at
$d/\lam \approx 10^{-13}$ --- sits just $1.2\times$ above the \texttt{snap} merge floor of the
closed-form spectrum. \textbf{That reasoning was wrong: family B does not use snapped truth.} It
takes the pair values from the mode formula directly, precisely so that a controlled separation
is not at the mercy of a merge policy. Measured, the two values are 833 floating-point ulps
apart at $\delta = 10^{-13}$ and 250 apart at $3 \times 10^{-14}$, and the pencil has
$d/(\varepsilon/c) = 138$ of headroom at the first.

Two such cases were therefore built and run, and both pass the floor precondition (headroom
$4.1$ and $3.3$). At $\pr = 10^{-4}$, $10^{-8}$ and $10^{-11}$ the pair may legally merge and
most methods pass cheaply; at $\pr = 10^{-14}$ it must be split, and \textbf{every method
fails}. So $\pr = 10^{-14}$ is now exercised as a genuine requirement, and the answer is
consistent with the ladder: nothing splits a pair at $d/\lam = 1.2\mathrm{e}{-13}$ when nothing
splits one at $3.6\mathrm{e}{-5}$.

\item \textbf{Rectangles only.} The property rectangles cannot reproduce is the tension floor
jumping between adjacent modes on a corner domain --- separately found to be largely an
artifact of a marginal basis, but re-measured at recipe grade on only one domain.

\item \textbf{Bracketing only.} A method that deliberately blurs bracketing into multiplicity
estimation is under-credited by construction. The swing is such a method: the same eigenvector
signal that cuts its cells also distinguishes a true degeneracy from two distinct eigenvalues,
and none of that is visible to a bracket-set metric.

\item \textbf{One budget, one seed, three settings per knob.} A method rejected here is
rejected \emph{at 1500 evaluations}, not in principle. The one non-monotonicity in the ladder
--- \texttt{lappy\_repaired} passing $d = 1.8\mathrm{e}{-2}$ while both swing variants pass the
tighter $5.9\mathrm{e}{-3}$ and fail it --- is most likely a knob-grid artifact, since
\texttt{per\_gap} was swept only over $\{6, 32, 160\}$, but that was not confirmed.

\item \textbf{The oracle column is best-of-$K$} and biased upward in the number of knobs. Only
the autonomous column may be compared across methods, and the gap between the two is the most
actionable quantity in the report.
\end{itemize}
